% Part: many-valued-logic
% Chapter: syntax-and-semantics
% Section: introduction

\documentclass[../../../include/open-logic-section]{subfiles}

\begin{document}

\olfileid[hi]{mvl}{syn}{int}

\olsection{परिचय}

चिरसम्मत तर्कशास्त्र में हम उन !!{formula} से काम करते हैं जो
प्रतिज्ञप्तिपरक संयोजकों $\lnot$, $\land$, $\lor$, $\lif$ और $\liff$ द्वारा
!!{propositional variable} से बनाए जाते हैं। चिरसम्मत तर्कशास्त्र का
अर्थविज्ञान परिभाषित करते समय हम दो सत्य-मान $\True$ और $\False$ उपयोग
करते हैं। हम !!{propositional variable} की व्याख्या किसी
!!a{valuation}~$\pAssign{v}$ में करते हैं, जो इन चरों को सत्य-मान $\True$,
$\False$ नियत करता है। तब कोई भी !!{valuation}, प्रत्येक !!{formula}~$!A$
का सत्य-मान $\pValue{v}(!A)$ निर्धारित करता है, और !!^a{formula} किसी
!!a{valuation}~$\pAssign{v}$ में संतुष्ट है, $\pSat{v}{!A}$, तभी और केवल
तभी जब $\pValue{v}(!A) = \True$।

बहुमूल्य तर्कशास्त्र, केवल $\True$ और $\False$ से अधिक सत्य-मानों की अनुमति
देकर चिरसम्मत द्विमूल्य तर्कशास्त्र का व्यापकीकरण करते हैं। अतः बहुमूल्य
तर्कशास्त्र में !!a{valuation}~$\pAssign{v}$ ऐसा फलन है जो प्रत्येक
!!{propositional variable}~$p$ को संभव सत्य-मानों की किसी श्रेणी का एक मान
नियत करता है। सामान्यतः हम अनुमत सत्य-मानों के समुच्चय को~$V$ कहेंगे।
चिरसम्मत तर्कशास्त्र ऐसा बहुमूल्य तर्कशास्त्र है जिसमें $V = \{\True,
\False\}$; और सत्य-मान~$\pValue{v}(!A)$ की संगणना संयोजकों की परिचित
अभिलाक्षणिक सत्य-सारणियों से होती है।

अतिरिक्त सत्य-मान जोड़ने पर ``और'', ``या'', ``नहीं'' और ``यदि---तो'' पढ़े
जाने वाले संयोजकों के लिए~$\pValue{v}(!A)$ संगणित करने का एक से अधिक
स्वाभाविक विकल्प होता है। इसलिए कोई बहुमूल्य तर्कशास्त्र केवल सत्य-मानों के
समुच्चय से नहीं, बल्कि प्रत्येक संयोजक के लिए चुने गए \emph{सत्य-फलनों} से
भी निर्धारित होता है। किंतु किसी बहुमूल्य तर्कशास्त्र~$\Log L$ के लिए इनके
चुन लिए जाने पर, चिरसम्मत तर्कशास्त्र की तरह सत्य-मान
$\pValue{v}(!A)[\Log L]$ मूल्यांकन से अद्वितीय रूप से निर्धारित होता है।
चिरसम्मत तर्कशास्त्र की तरह बहुमूल्य तर्कशास्त्र भी \emph{सत्य-फलनीय} हैं।

इन अर्थवैज्ञानिक आधार-घटकों से हम सर्वसत्यता, अनुगमन और संतुष्ट्यता की
अर्थवैज्ञानिक संकल्पनाओं के अनुरूप रूप परिभाषित कर सकते हैं। चिरसम्मत
तर्कशास्त्र में कोई !!a{formula} सर्वसत्य है यदि प्रत्येक~$\pAssign{v}$ के
लिए उसका सत्य-मान $\pValue{v}(!A) = \True$ हो। बहुमूल्य तर्कशास्त्र में इसे
भी कुछ व्यापक करना पड़ता है। पहली बात, यह आवश्यक नहीं कि सत्य-मानों के
समुच्चय~$V$ में~$\True$ हो। उदाहरणार्थ, कुछ बहुमूल्य तर्कशास्त्र $0$ और~$1$
के बीच की सभी परिमेय संख्याओं जैसे अंकों को अपने सत्य-मानों का समुच्चय
बनाते हैं। ऐसी स्थिति में $1$ सामान्यतः~$\True$ की भूमिका निभाता है। अन्य
तर्कशास्त्रों में केवल एक नहीं, अनेक सत्य-मान यह भूमिका निभाते हैं। अतः हम
माँग करते हैं कि प्रत्येक बहुमूल्य तर्कशास्त्र में \emph{निर्दिष्ट मानों}
का समुच्चय~$V^+$ हो। तब हम कह सकते हैं कि !!a{formula} किसी
!!a{valuation}~$\pAssign{v}$ में संतुष्ट है, $\pSat{v}{!A}[\Log L]$, तभी और
केवल तभी जब $\pValue{v}(!A)[\Log L] \in V^+$। !!^a{formula}~$!A$ उस
तर्कशास्त्र की सर्वसत्यता है, $\Entails[\Log L] !A$, तभी और केवल तभी जब
प्रत्येक~$\pAssign{v}$ के लिए $\pValue{v}(!A) \in V^+$। अंततः, हम कहते हैं
कि !!{formula} के किसी समुच्चय द्वारा $!A$ का अनुगमन होता है, $\Gamma
\Entails[\Log L] !A$, यदि~$\Gamma$ के सभी !!{formula} को संतुष्ट करने वाला
प्रत्येक !!{valuation}, ~$!A$ को भी संतुष्ट करे।

\end{document}
